\section*{Chapter \ref{ch:m2:how-bonds-trade} --- How Bonds Trade}

\begin{solution}{exo:m2:how-bonds-trade:1}
Only that it was larger than USD~5 million, at 99.50: the cap hides the exact size of
large investment-grade trades.
\end{solution}

\begin{solution}{exo:m2:how-bonds-trade:2}
$10 - 5e_n$: 10.00 basis points for one dealer, 7.18 for two, 4.85 for four. Without
leakage more is always better.
\end{solution}

\begin{solution}{exo:m2:how-bonds-trade:3}
Because the price reported for each bond in a basket is an allocation of one agreed
price for the whole basket, not a price negotiated for that bond; readers of the tape
must know not to treat it as one.
\end{solution}

\begin{solution}{exo:m2:how-bonds-trade:4}
6, 3 and 2 dealers.
\end{solution}

\begin{solution}{exo:m2:how-bonds-trade:5}
It buys ETF shares cheaply, redeems them for bonds and sells the bonds, earning the
discount and narrowing it. It may hesitate because the bonds' published values are
stale and it may not be able to sell them near those values, because redemption
delivers bonds it must then hold or fund, and because its own balance sheet is
constrained in stress.
\end{solution}

\begin{solution}{exo:m2:how-bonds-trade:6}
It may have no interest in the bond, no inventory or risk capacity, the client may
be one whose requests often precede price moves, the size may be too large to hedge,
or it may expect to lose the auction anyway and gain nothing but information.
\end{solution}

\begin{solution}{exo:m2:how-bonds-trade:7}
99.449: with a ten-second half-life the freshest quotes, 99.47 and 99.42, dominate and
the older ones count far less; the outlier is still dropped.
\end{solution}

\begin{solution}{exo:m2:how-bonds-trade:8}
Each dealer asked learns the client's intention; the losers can trade ahead of the
winner or mark the bond down, and the client pays for that leakage. The gain from an
extra bid shrinks as dealers are added, while the leakage cost grows with each, so
beyond a point more dealers raise the expected cost.
\end{solution}

\begin{solution}{pb:m2:how-bonds-trade:1}
\textbf{1.} 10.0 basis points, USD~25\,000.
\textbf{2.} 8.18, 7.77, 7.85 and 8.66 basis points.
\textbf{3.} Three dealers, 7.77 basis points, about USD~19\,400.
\textbf{4.} About USD~2\,200.
\textbf{5.} The expected best of $n$ bids grows ever more slowly with $n$.
\textbf{6.} Losing dealers who know a seller is active sell ahead, stop buying, or
mark the bond and its hedges down, so the price at which the manager trades, or its
next trade, is worse.
\textbf{7.} 6 dealers at 0.5 basis points, 2 at 1.5.
\textbf{8.} A large sale takes longer to absorb and moves the price more, so what the
losers learn is worth more to them.
\textbf{9.} Asking the dealers that are axed to buy, who are likely to bid best,
lets the manager ask fewer and leak less.
\textbf{10.} It adds bidders who have no dealer's incentive to trade ahead and may
want the bond; it may also leak to more parties, depending on the platform's design.
\textbf{11.} When selling many bonds at once, especially small or illiquid ones, it can
ask one price for the list and pay less per bond than in many RFQs.
\textbf{12.} The price and a size shown as ``5MM$+$'' if it is investment grade, within
fifteen minutes.
\textbf{13.} Buyers see that a large block traded but not its size, so they do not
know how much the dealer has left to sell.
\textbf{14.} If it holds ETF shares or can short them; selling ETF shares is not
selling the bond, but it may hedge the bond's credit exposure until the bond is sold.
\textbf{15.} Where the bond should trade, before asking, so it can judge the bids it
receives.
\textbf{16.} They price the client, the size and their own inventory; for a client whose
trades precede moves, they bid lower, and they may bid only to learn.
\textbf{17.} Compare prices after each RFQ with prices after comparable trades done
with fewer dealers, controlling for size and market moves.
\textbf{18.} High-yield bonds are less liquid and leakage is costlier, so fewer dealers
may be best; but private terms may be more dispersed, favouring more.
\textbf{19.} Named result: \emph{the best number of dealers} is three, at an expected
cost of 7.77 basis points, about USD~19\,400 on USD~25 million.
\textbf{20.} Because every dealer asked also learns what the client wants to do.
\end{solution}

\begin{solution}{iq:m2:how-bonds-trade:1}
A stock trades on exchanges with continuous public order books; a corporate bond,
one of many issues of the company, trades over the counter, mostly by requests for
quotes to dealers, with trades reported to TRACE after the fact. Liquidity is
concentrated in recent large issues; most bonds trade rarely.
\iqlookfor{OTC and RFQ, fragmentation by issue, post-trade transparency.}
\end{solution}

\begin{solution}{iq:m2:how-bonds-trade:2}
FINRA's system that collects reports of every OTC trade in eligible bonds within
fifteen minutes and publishes them. It shows price, time and size, but sizes of large
trades are capped, it shows no quotes or depth, and it shows trades only after they
happen.
\iqlookfor{what is published and what is hidden.}
\end{solution}

\begin{solution}{iq:m2:how-bonds-trade:3}
Model bids as a common value less a markup plus independent private terms; the best
of $n$ improves by $\sigma e_n$; each loser costs leakage $\ell$; minimise
$m - \sigma e_n + \ell(n-1)$. The answer rises with the dispersion of dealers' terms
and falls with leakage.
\iqlookfor{the trade-off and its comparative statics.}
\end{solution}

\begin{solution}{iq:m2:how-bonds-trade:4}
Value each bond from composites and the issuer curves; aggregate the basket's risk
(duration, spread duration, sectors, liquidity); price the cost of hedging it with
index swaps and ETFs and of working out the bonds over days, bond by bond according
to liquidity; add a charge for the basket's residual risk and capital; quote one
price.
\iqlookfor{basket-level risk, hedging instruments and liquidation cost.}
\end{solution}

\begin{solution}{iq:m2:how-bonds-trade:5}
Its net asset value uses bond prices that may be stale or estimated, while the ETF
trades continuously; in stress, sellers use the ETF and its price falls first. The
authorised participants' arbitrage is limited by the cost and risk of trading the
bonds. The ETF price is often the better measure of where the bonds could actually be
sold.
\iqlookfor{stale NAV and limits to arbitrage.}
\end{solution}

\begin{solution}{iq:m2:how-bonds-trade:6}
Precompute composites and hedging sensitivities for every bond and update them on
market events; on each RFQ, look up the bond, apply client, size and inventory
adjustments from cached tables, check limits, and answer; keep the path free of I/O
and allocation; log decisions for later analysis; monitor hit ratios and mark-outs to
recalibrate.
\iqlookfor{precomputation, fast lookups, limits, feedback.}
\end{solution}
